\DTMsavedate{mydate}{2022-06-30}
\maketitle{Graphics for Electronic Engineers}{Electronic Engineering}{\DTMusedate{mydate}}

% ----------------------------------------------------- %
% COURSE INFO
\subsection*{Course Information}
\vspace{0ex}%
\noindent\parbox{0.5\textwidth}{%
    \noindent\begin{tabular}{@{}ll}
                 \textsf{Course:}          & 	Graphics for Electronic Engineers \\
                 \textsf{Course Code:}         & 	TEE 1201 \\
                 \textsf{Credit Hours:}    & 	10
    \end{tabular}}
\vspace{2ex}\hrule\vspace{2ex}

% ----------------------------------------------------- %
% INSTRUCTOR INFO
\subsection*{Lecturer Information}
\noindent\parbox{0.5\textwidth}{%
    \noindent\begin{tabular}{@{}ll}
%             \input{\staff/komichi.tex}
    \end{tabular}}
\vspace{2ex}\hrule\vspace{2ex}


% ----------------------------------------------------- %
% COURSE SYNOPSISDESCRIPTION
\subsection*{Course Synopsis}
The module focuses on graphical techniques for drawing circuit diagrams, logic circuits, flowcharts.
It covers concepts of engineering drawings, presentation of graphs and design of artwork for printed circuit boards.
Use of pictures and cartoons and use of CIRCUIT MAKER PRO program for graphical design.
\vspace{2ex} \hrule\vspace{2ex}

% ----------------------------------------------------- %
% LEARNING OBJECTIVES
\subsection*{Learning Objectives}
On completion of the course students should be able to:
\begin{outline}
    \1

    \1

    \1

    \1

    \1

    \1

    \1

    \1

    \1

    \1

    \1

\end{outline}
\vspace{2ex}\hrule\vspace{2ex}


% ----------------------------------------------------- %
% COURSE TOPICS
\subsection*{Topics}
\begin{outline}
    \1
    \begin{itemize}
        \item
        \item
        \item
        \item
        \item
    \end{itemize}
    \1
    \1
\end{outline}
\vspace{2ex} \hrule\vspace{2ex}

% ----------------------------------------------------- %
% Students Assenssment
\subsection*{Assessment of Students}
\begin{outline}
    \1 Students are assessed through:
    \begin{enumerate}
        \item course work consisting of laboratory exercises.
        \item end of semester final laboratory examination.
    \end{enumerate}

\end{outline}
\vspace{2ex}\hrule\vspace{2ex}

% ----------------------------------------------------- %
% Grade Evaluation
\subsection*{Course Evaluation and Grading Policies}
Grades are based on:
Continuous assessment (\qty{25}{\percent}),
Final Exam (\qty{75}{\percent})
\vspace{2ex}\hrule
\vspace{2ex}\hrule\vspace{2ex}